% Part: turing-machines
% Chapter: undecidability
% Section: enumerating-tms

\documentclass[../../../include/open-logic-section]{subfiles}

\begin{document}

\olfileid[hi]{tur}{und}{enu}
\olsection{ट्यूरिंग मशीनों का परिगणन}

\begin{explain}
हम दिखा सकते हैं कि सभी ट्यूरिंग मशीनों का समुच्चय !!{enumerable} है।
यह इस तथ्य से निकलता है कि प्रत्येक ट्यूरिंग मशीन का परिमित वर्णन किया
जा सकता है। अवस्थाओं का समुच्चय और पट्टी की शब्दावली परिमित समुच्चय हैं।
संक्रमण फलन $Q \times \Sigma$ से
$Q \times \Sigma \times \{\TMleft, \TMright, \TMstay\}$ तक का आंशिक
फलन है; अतः जिन परिमित संख्या में युग्म-आगतों पर वह परिभाषित है, उनके
लिए उसके मान सूचीबद्ध करके उसे भी निर्दिष्ट किया जा सकता है।

इतना कथन सही है, किंतु यहाँ एक सूक्ष्म अंतर है। ट्यूरिंग मशीन की परिभाषा
ने इस पर कोई प्रतिबंध नहीं लगाया कि अवस्थाओं के समुच्चय और पट्टी-वर्णमाला
में कौन-से !!{element} हो सकते हैं। अतः उदाहरण के लिए, तकनीकी रूप से
प्रत्येक वास्तविक संख्या के लिए ऐसी ट्यूरिंग मशीन है जो उस संख्या को एक
अवस्था के रूप में उपयोग करती है। फिर भी ट्यूरिंग मशीन का \emph{व्यवहार}
इससे स्वतंत्र है कि कौन-सी वस्तुएँ अवस्थाओं और शब्दावली का काम करती हैं।
\olref{fig:variants} की दोनों ट्यूरिंग मशीनों पर विचार करें।
\begin{figure}
\begin{center}
\begin{tikzpicture}[->,>=stealth',shorten >=1pt,auto,node distance=2.8cm,
                    semithick]
  \tikzstyle{every state}=[fill=none,draw=black,text=black]

  \node[initial,state]         (A)              {$q_0$};
  \node[state]         (B) [right of=A] {$q_1$};

  \path (A) edge [bend left] node {\TMtrans{\TMstroke}{\TMstroke}{\TMright}} (B)
        (B) edge [loop above] node {\TMtrans{\TMblank}{\TMblank}{\TMright}} (B)
            edge [bend left] node {\TMtrans{\TMstroke}{\TMstroke}{\TMright}} (A);
\end{tikzpicture}\\
\begin{tikzpicture}[->,>=stealth',shorten >=1pt,auto,node distance=2.8cm,
  semithick]
\tikzstyle{every state}=[fill=none,draw=black,text=black]

\node[initial,state]         (A)              {$s$};
\node[state]         (B) [right of=A] {$h$};

\path (A) edge [bend left] node {\TMtrans{A}{A}{\TMright}} (B)
(B) edge [loop above] node {\TMtrans{\TMblank}{\TMblank}{\TMright}} (B)
edge [bend left] node {\TMtrans{A}{A}{\TMright}} (A);
\end{tikzpicture}
\end{center}
\caption{\emph{सम} मशीन के रूपांतर}
\ollabel{fig:variants}
\end{figure}
ये दोनों आरेख दो मशीनों के अनुरूप हैं: $M$, जिसकी पट्टी-वर्णमाला
$\Sigma = \{\TMendtape,\TMblank,\TMstroke\}$ और अवस्था-समुच्चय
$\{q_0,q_1\}$ हैं; तथा $M'$, जिसकी वर्णमाला $\Sigma' =
\{\TMendtape,\TMblank,A\}$ और अवस्थाएँ $\{s,h\}$ हैं।
किंतु अन्यथा उनके निर्देश समान हैं: $M$, $n$ संख्या के $\TMstroke$ के
अनुक्रम पर तभी रुकती है जब $n$ सम हो, और $M'$, $n$ संख्या के $A$ के
अनुक्रम पर तभी रुकती है जब $n$ सम हो। हमने केवल $\TMstroke$ का नाम
बदलकर~$A$, $q_0$ का~$s$, और $q_1$ का~$h$ किया है। इस उदाहरण का
सामान्यीकरण होता है: जब तक एक ट्यूरिंग मशीन दूसरी से प्रतीकों और अवस्थाओं
के ऐसे पुनर्नामकरण द्वारा मिलती है, हम दोनों को एक ही मान सकते हैं।
वस्तुतः हम ट्यूरिंग मशीन के प्रतीकों और अवस्थाओं को केवल धन पूर्णांक मान
सकते हैं: $\sigma_0$ के स्थान पर~$1$, $\sigma_1$ के स्थान पर~$2$,
इत्यादि; $\TMendtape$ को~$1$, $\TMblank$ को~$2$, इत्यादि। इस प्रकार
\emph{सम} मशीन \olref{fig:standard-even} में दर्शाई मशीन बन जाती है।
\begin{figure}
\[\begin{tikzpicture}[->,>=stealth',shorten >=1pt,auto,node distance=2.8cm,
  semithick]
\tikzstyle{every state}=[fill=none,draw=black,text=black]

\node[initial,state]         (A)              {$1$};
\node[state]         (B) [right of=A] {$2$};

\path (A) edge [bend left] node {\TMtrans{3}{3}{\TMright}} (B)
(B) edge [loop above] node {\TMtrans{2}{2}{\TMright}} (B)
edge [bend left] node {\TMtrans{3}{3}{\TMright}} (A);
\end{tikzpicture}
\]
\caption{एक मानक \emph{सम} मशीन}
\ollabel{fig:standard-even}
\end{figure}
जिस ट्यूरिंग मशीन की अवस्थाएँ और प्रतीक धन पूर्णांक हों, उसे हम
\emph{मानक} मशीन कह सकते हैं और अब से केवल मानक मशीनों पर विचार कर
सकते हैं।\footnote{``मानक मशीन'' पदावली स्वयं मानक नहीं है।}

हम दिखाना चाहते थे कि ट्यूरिंग मशीनों का समुच्चय !!{enumerable} है, और
ऊपर के विचारों को ध्यान में रखते हुए मानक ट्यूरिंग मशीनों के समुच्चय को
!!{enumerable} दिखाना पर्याप्त है। मान लें कि हमें मानक ट्यूरिंग मशीन
$M = \tuple{Q, \Sigma, q_0, \delta}$ दी गई है। धन पूर्णांकों की परिमित
अक्षर-शृंखला से उसका वर्णन कैसे करें? पहले हम अवस्थाओं की संख्या, स्वयं
अवस्थाएँ, प्रतीकों की संख्या, स्वयं प्रतीक और आरंभिक अवस्था सूचीबद्ध
करेंगे। (याद रखें, ये सभी धन पूर्णांक हैं क्योंकि $M$ एक मानक मशीन है।)
अब~$\delta$ का क्या करें? संभावित आगतों का समुच्चय, अर्थात् युग्मों
$\tuple{q,\sigma}$ का समुच्चय, परिमित है क्योंकि $Q$ और~$\Sigma$
परिमित हैं। अतः~$\delta$ में निहित सूचना उन सभी $5$-गुणजों
$\tuple{q, \sigma, q', \sigma', d}$ की परिमित सूची मात्र है जिनके लिए
$\delta(q,\sigma) = \tuple{q', \sigma', D}$ है और $d$, दिशा~$D$ को
कूटबद्ध करने वाली संख्या है (मान लें, $1$ से~$\TMleft$, $2$
से~$\TMright$, और $3$ से~$\TMstay$ कूटबद्ध होते हैं)।

इस प्रकार प्रत्येक मानक ट्यूरिंग मशीन का वर्णन धन पूर्णांकों की एक
परिमित सूची द्वारा, अर्थात् अनुक्रम $s_M \in
(\PosInt)^*$ के रूप में,
किया जा सकता है। उदाहरणतः मानक \emph{सम} मशीन को निम्न अनुक्रम कूटबद्ध
करता है:
\[
2, \underbrace{1, 2}_Q, 3, \overbrace{1, 2, 3}^\Sigma, 1, \underbrace{1, 3, 2, 3, 2}_{\delta(1,3) = \tuple{2,3,R}}, 
\overbrace{2, 2, 2, 2, 2}^{\delta(2,2) = \tuple{2,2,R}},
\underbrace{2, 3, 1, 3, 2}_{\delta(2,3) = \tuple{1,3,R}}.
\]
\end{explain}

\begin{thm}
$\Nat$ से~$\Nat$ तक ऐसे फलन हैं जो ट्यूरिंग-संगणनीय नहीं हैं।
\end{thm}

\begin{proof}
हम जानते हैं कि धन पूर्णांकों के परिमित अनुक्रमों का समुच्चय
~$(\PosInt)^*$, !!{enumerable} है
(\cref{sfr:siz:zigzag:prob:posint-star})। इससे मिलता है कि
~$(\PosInt)^*$ के उपसमुच्चय के रूप में मानक ट्यूरिंग मशीनों के वर्णनों
का समुच्चय स्वयं परिगणनीय है। $\Nat$ से~$\Nat$ तक प्रत्येक
ट्यूरिंग-संगणनीय फलन की संगणना कोई (वस्तुतः अनेक) ट्यूरिंग मशीन करती
है। उसकी अवस्थाओं और प्रतीकों का नाम बदलकर धन पूर्णांक रखने पर (विशेषतः
$\TMendtape$ को~$1$, $\TMblank$ को~$2$, और $\TMstroke$ को~$3$) हम देख
सकते हैं कि प्रत्येक ट्यूरिंग-संगणनीय फलन की संगणना कोई मानक ट्यूरिंग
मशीन करती है। इसका अर्थ है कि $\Nat$ से~$\Nat$ तक के सभी
ट्यूरिंग-संगणनीय फलनों का समुच्चय भी परिगणनीय है।

दूसरी ओर, $\Nat$ से~$\Nat$ तक के सभी फलनों का समुच्चय !!{enumerable}
नहीं है (\cref{sfr:siz:red:prob:nat-nat})। यदि सभी फलन किसी ट्यूरिंग
मशीन द्वारा संगणनीय होते, तो उनकी संगणना करने वाली ट्यूरिंग मशीनों के
सभी वर्णनों को सूचीबद्ध करके हम सभी फलनों के समुच्चय का परिगणन कर सकते
थे। अतः कुछ फलन ट्यूरिंग-संगणनीय नहीं हैं।
\end{proof}

\begin{prob}
  क्या आप ट्यूरिंग मशीनों का वर्णन करने की ऐसी रीति सोच सकते हैं जिसमें
  अवस्थाओं और वर्णमाला के प्रतीकों को स्पष्टतः सूचीबद्ध करना आवश्यक न
  हो? आप ``मानक'' मशीन की अपनी धारणा परिभाषित कर सकते हैं, किंतु यह भी
  बताएँ कि आपके नए अर्थ की किसी ``मानक'' मशीन द्वारा प्रत्येक ट्यूरिंग
  मशीन का अनुकरण क्यों किया जा सकता है।
\end{prob}

\end{document}
